\section{Counting frames and fixing images}\label{sec:frame-laws}

The cost of prescribing frame images depends on the number of independent
directions observed. We count their possible images first.

\subsection{Quadratic Gram data}

For vectors \(v_1,\ldots,v_a\), their quadratic Gram data consist of
\(q_0(v_i)\) and \(\beta_0(v_i,v_j)\), for \(i<j\).
There are \(a(a+1)/2\) entries. Each diagonal entry records a quadratic value.

% lax begin Lax793645.QuadraticFrameCounts
\begin{lemma}\label{lem:quadratic-frame-count}
Suppose \(v_1,\ldots,v_a\) are independent in a hyperbolic space of
full dimension \(N\). For any prescribed \(b_1,\ldots,b_a,q\), the
number of vectors \(v\) satisfying
\[
 \beta_0(v_i,v)=b_i,\qquad q_0(v)=q
\]
is \(2^{N-a-1}(1+O(2^{a-N/2}))\). The number of such vectors outside
the previous span depends only on the quadratic Gram data and the
prescribed entries. Consequently every partial isometric embedding has
the same number of completions.
\end{lemma}
% lax end
% lax begin Lax793645Proofs.Manuscript.QuadraticFrameCounts.quadratic_column_mass
\begin{proof}
Nondegeneracy makes the polar equations independent. Their solution set
has size \(2^{N-a}\). The quadratic Gauss sum on that set is
\[
 2^{N/2-a}\sum_{t\in\F^a}
 (-1)^{q_0(\sum_i t_iv_i)+t\cdot b}.
\]
To obtain this identity, expand each polar equation as
\(\frac12\sum_{t_i\in\F}(-1)^{t_i(\beta_0(v_i,v)+b_i)}\),
then translate \(v\) by \(\sum_i t_iv_i\) in the hyperbolic Gauss sum.
The sum has absolute value at most \(2^{N/2}\). It is the number of
solutions with quadratic value 0 minus the number with value 1. Their
sum is \(2^{N-a}\), which gives the asserted count for either value.

A vector in the old span is \(\sum_i t_iv_i\). Whether it satisfies the
new equations is determined by the old Gram data. Thus the number excluded
by independence is determined by those data as well. In particular,
two lists with the same Gram data admit equally many choices for their
next independent vector. Complete the first list to a basis. At each step
there is a matching choice for the second list, since the first has at
least one choice. The linear map between the resulting bases preserves
the quadratic form. This uses the exact count at every dimension.
\end{proof}
% lax end

For \(a\le N/4\), it follows that a uniform list of ambient vectors has
prescribed independent quadratic Gram data with probability
\begin{equation}\label{eq:quadratic-gram-count}
 2^{-a(a+1)/2}\bigl(1+O(a2^{a-N/2})\bigr).
\end{equation}
The same count gives the frames in \eqref{eq:one-frame}. If
\(N\ge 2a+4\), at least \((2^{N-a}-2^{N/2})/2\) vectors realize
any prescribed next entries. This exceeds \(2^a\), the size of the
previous span. Appending independent columns therefore gives an embedding
with the required Gram data.

The frame law prescribes only part of the full Gram data. Sum the completion
counts over the unspecified entries. This shows that an independent list
of primal directions has exactly the uniform image law with its prescribed
quadratic Gram data. It also gives the following estimates, used repeatedly.
For every fixed \(\xi>0\), by choosing \(M_0\) large enough, any consistent
prescription of \(t\) new image directions has raw probability at most
\begin{equation}\label{eq:image-cost}
 2^{-(1-\xi)tN}.
\end{equation}
The bound is uniform over the already stored bounded set of directions.
For free entries of a bounded cross table, every consistent assignment
has probability \(2^{-c+o(N)}\), where \(c\) is the number of independent
scalar equations. A prescribed kernel must be counted separately from
these Gram equations. We do that explicitly in \Cref{sec:histograms}.

\subsection{Leaves with stored images}

A \emph{pin} specifies the exact images of a coefficient subspace under
the unit's frame maps. A \emph{leaf} is a conditional law together with
its pins. We require the pin space to be closed under projection onto the
primal and the two channel blocks at each of the two endpoints.
There are six projections.

For a leaf, a further exact-image event has \emph{new rank} \(t\) if it
prescribes \(t\) coefficient directions independent modulo the stored
ones. Call the leaf \(\zeta\)-spreading if every such event has probability
at most \(2^{-(1-\zeta)tN}\). Only the bounded ranks used in the subsequent
queries are needed.

% lax begin Lax793645.ClosedPins
\begin{lemma}[Closed pins]\label{lem:closed-peeling}
From a unit law with a fixed exponential density cap one can retain
at least half its mass as a mixture of spreading leaves. Their total pin
dimension on each leaf has a bound independent of \(n\). The pins
may be required to be closed under the six block projections. Further
restriction to a set of mass \(2^{-o(N)}\) allows another such decomposition,
with a larger fixed pin budget and twice the original spreading error.
\end{lemma}
% lax end
% lax begin Lax793645Proofs.Manuscript.ClosedPins.one_first_leaves
\begin{proof}
If a leaf violates spreading, condition on a violating exact-image event
of new rank \(t\). Close its coefficient space under the block projections.
Let \(u\) be the total new dimension after closure, including the
original \(t\) directions. Each direction has at most 6 projections, so
\(t\le u\le6t\). There are at most \(2^{(u-t)N}\) ways to specify the
remaining images. Some choice has conditional probability at least
\(2^{-(u-t)N}\). The whole step thus has mass at least
\[
 2^{-(1-\zeta)tN}2^{-(u-t)N}
 \ge2^{-(1-\zeta/6)uN}.
\]
The raw image bound \eqref{eq:image-cost}, with \(\xi<\zeta/6\), charges
strictly more for the same stored directions. After total new dimension
\(k\), the likelihood ratio has gained at least
\(2^{(\zeta/6-\xi)kN}\). The initial density cap bounds this ratio.
Choose the pin budget larger than the cap exponent divided by
\(\zeta/6-\xi\). No branch of this procedure can exceed that budget.

To form a mixture, keep a residual subset of the original finite law.
While it has mass at least \(1/2\), normalize it and run the preceding
procedure until a spreading leaf is obtained. Remove that leaf from the
residual set and record its original probability. The residual law has
at most twice the initial density cap, so one pin budget works at every
iteration. Each removal has positive mass. Finiteness makes the process
terminate, and the recorded leaves together retain at least \(1/2\) of
the law. Their weights are the actual removed masses.

For the second assertion, start with the old pins and allow a fixed reserve
of new directions. Restricting by mass \(2^{-o(N)}\) adds only \(o(N)\)
to the logarithm of the density cap. That cost fits in the reserve for
large \(n\). The same construction, with twice the spreading error,
gives the new leaves.
\end{proof}
% lax end

Let \(K\) bound the total pin dimension. Write \(P_{i,e}\le B\) for the
primal pins at endpoint \(i\), component \(e\). Closure under the block projections ensures that the images of these
primal directions are among the stored data.
For a subspace \(D\), write
\[
 D^{(2)}=\Span\{vv\trans:v\in D\}.
\]
The effective profiles at endpoint \(i\) form
\begin{equation}\label{eq:effective-space}
 C_i=\{x\in\cX:x_e\in P_{i,e}^{(2)}\text{ for every }e\}.
\end{equation}
They have dimension at most \(K^2\), and every such profile has total
component rank at most \(K\). Choose a basis on each leaf. Evaluations
on that basis will be part of the finite status of a query.

Later restrictions preserve the old pins and effective profiles. The
retained mass \(\theta_n\) is bounded below by a positive constant in
Proposition~\ref{prop:status-preparation}. We use this bound in the final probability
comparison.
